% SPDX-License-Identifier: LPPL-1.3c
% Copyright (C) 2026 Christophe Jorssen
% Author and Current Maintainer: Christophe Jorssen <christophe.jorssen@gmail.com>
% LPPL maintenance status: maintained
% This file is part of luafemm. See LICENSE and LICENSES.md for its terms.
\section{Boundary conditions}
\label{sec:boundaries}
The outer rectangle is part of the physical model. Leaving its four sides
at the default $A_z=0$ confines flux to a finite box. This is often useful
for a first calculation, but the boundary should be moved farther away
to assess its influence. Other conditions can impose a background field,
describe a symmetry plane, or specify a mixed relation between potential
and tangential field.

\subsection{A uniform field to start with}
Suppose that the empty domain should contain $B_x=0.2\,\mathrm T$ and
$B_y=0$. Since $B_x=\partial_y A_z$, choose $A_z=0.2y$, with $y$ in
metres. The rectangle below uses millimetres for geometry; the coefficient
is nevertheless entered as |.2|, because boundary data always use SI.

\begin{codeexample}[width=5cm]
\begin{tikzpicture}[femm/problem={
  xmin=-20,xmax=20,ymin=-15,ymax=15,
  mesh size=4},
  femm/boundary={type=dirichlet,
    y coefficient=.2}]
  \draw (-20,-15) rectangle (20,15);
  \pic[femm/lines=9] {femm field};
\end{tikzpicture}
\end{codeexample}

More generally, use |x coefficient=|$-B_y$ and
|y coefficient=|$B_x$ to prescribe an affine potential on the exterior.
Adding a material region changes the interior field. Imposing the
background potential on a finite rectangle still perturbs the exterior
solution: enlarge the domain and compare results.

\subsection{Command and key reference}
\begin{key}{/tikz/femm/boundary=\marg{options} (default \normalfont empty)}
Declare a boundary condition using ordinary TikZ keys. Its subkeys are
exactly those of |\femmboundary| below. In picture options, declarations
are collected and applied after the new model is created, before the
picture body starts. This key may precede or follow |femm/problem|.
Repeated declarations are applied in order; the last one for each side
wins. Setting |femm/problem| in several styles still creates only one
model, using the final problem options.

Normal TikZ |.style|, |.append style|, parameterised styles and
|every picture| are supported. Outside a picture, a direct
|\tikzset{femm/boundary=...}| records defaults for subsequent magnetic
pictures in that TeX scope. Each picture inherits its own copy; body
declarations do not change the defaults for later pictures. Pending
options, including macros and boundary-family styles, are evaluated
when the model starts. A picture without |femm/problem| does not apply
them to a previous model.

In the body of a started magnetic picture, the key applies immediately.
For example, |\tikzset{femm/boundary={side=left,type=neumann}}| is equivalent
to the corresponding command. It can also be used in |scope| options.
Physical declarations persist after the graphical scope ends: they are
not temporarily active boundary conditions. A declaration after meshing
raises an error, just as with the command.

This key selects complete sides of the model rectangle. On a drawing
path, it still modifies the selected side; it does not assign a boundary
condition to that path's contour. Prefer picture options or a dedicated
setup scope. Do not put it in |every path| or in a style used to draw
computed fields, since those paths may be drawn after the model is meshed.
Using the key without a value resets all sides to zero Dirichlet data.
\end{key}

\begin{codeexample}[code only]
\tikzset{
  half model/.style={
    femm/boundary={side=left,type=neumann}},
  background field/.style={
    femm/boundary={y coefficient=#1}},
  background field/.default=.2}
\begin{tikzpicture}[half model,
  femm/problem={xmin=0,xmax=85,ymin=-75,ymax=95,mesh size=4}]
  % Declare the half-device, then request a field or profile.
\end{tikzpicture}
\end{codeexample}

\begin{command}{\femmboundary\oarg{options}}
Declare a condition on one complete side of the rectangular domain, or
on all four sides. Call this command in the picture body, after the model
has been created and before the first mesh, solve or field request.
It produces no drawing. The side names refer to the original model frame;
a local |rotate| or |scale| does not change their meaning.

The keys belong to |/femm/boundary|. Each call first resets the six keys
to their defaults below, then processes its options. A later declaration
replaces the earlier condition on the selected side. Other sides keep
their conditions. Declarations persist in the current model even if the
command is inside a TeX group. They cannot modify a meshed model; start a
new picture for a different physical calculation.

All four numerical keys accept PGF mathematical expressions. Values are
expanded, evaluated with PGF's floating-point library, and passed to Lua
as finite numbers. No user text is evaluated as Lua code. Normal PGF
styles, including |.style| and |.append style|, are available in the family.
The command and |femm/boundary| share this implementation and can be
mixed in one model. Use the TikZ key in picture options; the command
itself requires an already created model.
\end{command}

\begin{key}{/femm/boundary/side=\meta{choice} (initially all)}
The |.is choice| alternatives are |all|, |left|, |right|, |bottom| and |top|.
They select $x=x_{\min}$, $x=x_{\max}$, $y=y_{\min}$ and $y=y_{\max}$,
respectively. A condition applies to the whole side, including material
regions which reach it. To configure several sides differently, repeat
the command. There is no automatic symmetry detection.
\end{key}
\begin{key}{/femm/boundary/type=\meta{choice} (initially dirichlet)}
A |.is choice| selector with alternatives |dirichlet|, |neumann| and
|robin|. Their physical meanings are given below. Every side initially
has homogeneous Dirichlet data, whether or not this command is used.
\end{key}
\begin{key}{/femm/boundary/value=\meta{expression} (initially 0)}
Constant term $g_0$ in $g(x,y)=g_0+g_xx+g_yy$, with coordinates in metres
relative to the model origin. Its unit is T\,m for Dirichlet data and
A/m for Neumann or Robin data.
\end{key}
\begin{keylist}[/femm/boundary]{x coefficient=\meta{expression} (initially 0),
  y coefficient=\meta{expression} (initially 0)}
The coefficients $g_x,g_y$ of the affine data. Their units are T for
Dirichlet and A/m$^2$ for Neumann or Robin. These are spatial coefficients,
not components of H; their effect also depends on the selected side.
The potential and field use the model's physical frame, independently of
graphic transformations.
\end{keylist}
\begin{key}{/femm/boundary/coefficient=\meta{nonnegative expression} (initially 0)}
The Robin coefficient $c$, in A/(T\,m$^2$). A nonzero value requires
|type=robin|. Negative coefficients are rejected because the present
solver requires a positive definite tangent after constraints are applied.
Zero makes a Robin condition identical to Neumann data.
\end{key}

\begin{codeexample}[code only]
\pgfkeys{/femm/boundary/symmetry/.style={type=neumann,value=0}}
% Inside a picture, before meshing:
\femmboundary[side=left,symmetry]
% All other sides retain their initial Az=0 condition.
\end{codeexample}

\subsection{Signs, corners and physical interpretation}
Let $\mathbf n=(n_x,n_y)$ be the outward unit normal and
$\mathbf t=(-n_y,n_x)$ the counterclockwise unit tangent. Write
$H_t=\mathbf H\cdot\mathbf t$, including the coercive contribution
$\mathbf H=\nu\mathbf B-H_c\mathbf m$ in a permanent magnet.
The three conditions are
\[
 \begin{array}{ll}
 \texttt{dirichlet}:&A_z=g,\\
 \texttt{neumann}:&H_t=g,\\
 \texttt{robin}:&H_t=cA_z+g.
 \end{array}
\]
On the left, right, bottom and top sides, $H_t$ is respectively
$-H_y$, $H_y$, $H_x$ and $-H_x$. The sign does not depend on how a
decorative rectangle is drawn. This orientation also gives Ampere's law
$\oint H_t\,\mathrm ds=\int J_z\,\mathrm dS$.

A constant Dirichlet value makes the normal component of B vanish on
that side: field lines run along it. Homogeneous Neumann data make the
tangential component of H vanish. In air and an isotropic unmagnetised
material, B is then normal to the side, which is useful for the appropriate
symmetry plane. In a permanent magnet, $H_t=0$ does not generally imply
$B_t=0$. The condition concerns the total H, not just a derivative of A.

At a corner shared by Dirichlet and weak data, the prescribed potential
wins at that node; the weak edge integral still contributes to its free
neighbours. Two Dirichlet sides must prescribe the same corner value.
Conflicts raise an error rather than silently choosing a side. Affine
values are compared with a relative tolerance of $10^{-10}$ and a
$10^{-12}$ T\,m comparison scale floor.

In an unmagnetised material, the Robin relation can also be written
\[
 \nu\,\partial_n A_z+cA_z+g=0.
\]
This corresponds to FEMM's static mixed condition with $c_0=c$ and
$c_1=g$. A positive Robin coefficient is not, by itself, an exact
open-boundary treatment. In particular, a circular asymptotic formula
cannot simply be transferred to this rectangular domain without an
approximation study.

\subsection{Solving only a symmetric half}
The complete example \texttt{examples/boundary-conditions.tex} models the
right half of a U electromagnet with mirrored windings on both legs.
The cut is $x=0$, where H has no tangential component. The only change to
the boundary setup is
\begin{codeexample}[code only]
\begin{tikzpicture}[femm/problem={
  xmin=0,xmax=85,ymin=-75,ymax=95,mesh size=4},
  femm/boundary={side=left,type=neumann}]
  % Declare the right half of the core, armature and winding.
  % Refine the right gap, then draw the field and register a profile.
\end{tikzpicture}
\end{codeexample}
The other sides remain at $A_z=0$. The example plots the gap field with
PGFPlots. The symmetry assumption concerns both geometry and excitation;
the original tutorial's single winding on one leg does not have this
exact mirror symmetry. A normal material interface needs no boundary
declaration: continuity is already part of the finite-element formulation.

The complete source and computed result follow. The style declares the
symmetry condition before |femm/problem| in the option list; it is applied
after model creation. The field lines cross the symmetry cut normally,
and the graph samples the remaining air gap.
\luafemmexample{boundary-conditions}

\subsection{When no potential is prescribed}
With only Neumann conditions (including Robin with $c=0$), A is determined
only up to a constant. The package first checks the discrete compatibility
condition $\int J_z\,\mathrm dS=\oint g\,\mathrm ds$, with a relative
tolerance of $10^{-10}$ against the sum of absolute load contributions.
It then fixes $A_z=0$ at the first mesh vertex, the lower-left corner for
the supported meshers. This is a gauge reference, not a magnetic boundary.
Incompatible data raise an error; a reference node must not conceal a
current imbalance. Any Dirichlet side or positive Robin coefficient
removes this constant nullspace and suppresses the extra reference.

Boundary types and numerical values are included in the solution cache
descriptor. Changing them reuses a matching mesh and solves again;
|mode=frozen| rejects an outdated solution. The cache revision was
increased for this feature; earlier caches are rebuilt in |auto| mode.

This release supports complete sides of the rectangular exterior, affine
data and constant nonnegative Robin coefficients. Boundary conditions on
arbitrary internal curves, partial sides, periodic/antiperiodic pairs,
sliding air gaps, circular infinite elements and automatic open-boundary
constructions remain future work. Domains with unmeshed holes are a
separate topology extension.
